\begin{abstract}
We ask whether induction head circuits in attention-only transformers remain
structurally stable under fine-tuning on a narrow distribution, and at what
training step any structural change first occurs.
%
We fine-tune the two-layer \texttt{attn-only-2l} transformer on 500{,}000
tokens of Python code and a matched TinyStories prose control, measuring
per-head induction scores and activation-patching attribution at every
100-step checkpoint; at initialisation, L1H6 has a prefix-matching score
of $0.408 \pm 0.103$ (12.4$\times$ the next head) and attribution $0.952$,
replicating the canonical induction circuit of \citet{olsson2022context}.
%
We find that the induction circuit \emph{strengthens} rather than degrades
under both fine-tuning conditions: L1H6's prefix-matching score increases
from $0.408$ to $0.584$ (code) and $0.583$ (prose) after 100 steps, and
to $0.638$ and $0.616$ after 200 steps, with the largest change in the
first interval and near-identical trajectories across both domains
--- a result not previously reported for this model class ---
implying that, at a 500K-token budget, the induction circuit is not silently
eroded by fine-tuning and may instead be reinforced irrespective of domain.
%
These results demonstrate that circuit-level diagnostics are a necessary
complement to capability benchmarks when evaluating safety-relevant
properties of fine-tuned language models, not only to detect degradation
but also to detect unexpected reinforcement.
\end{abstract}
